
% Introdução 
\chapter[Introdução]{Introdução}\label{capitulo1}
\addcontentsline{toc}{chapter}{Introdução}

% Marcador visível para pendências de fonte (remover antes da versão final).
\providecommand{\completar}[1]{\textcolor{red}{[COMPLETAR: #1]}}

\section{Considerações Iniciais}

O transporte público é uma alternativa de mobilidade que pode reduzir os congestionamentos e a poluição do ar, além de favorecer sociedades mais sustentáveis \cite{hussain2017}. Com a expansão da internet e das tecnologias móveis, surgiram aplicativos que auxiliam o passageiro a planejar suas viagens, fornecendo informações em tempo real e serviços como a venda de bilhetes \cite{hussain2017}. Nesse cenário, a interface do aplicativo influencia a adoção e a satisfação dos usuários \cite{suzianti2020}.

Em Divinópolis, Minas Gerais, o aplicativo FaleBus permite ao usuário controlar o cartão DivPass e fazer recargas por cartão de crédito ou por PIX. De acordo com os dados de divulgação do aplicativo, ele é oferecido pela TransCid, foi lançado em março de 2017 e, na versão 4.3.7, atualizada em dezembro de 2025, tinha mais de 100 mil downloads. \completar{citar a fonte desses dados e a data de consulta}

Aplicativos como esse concentram funções que o passageiro usa com frequência, e dificuldades de navegação, de compreensão das informações ou de execução de tarefas podem prejudicar sua experiência. Antes de propor mudanças, porém, é preciso saber quais dificuldades de fato ocorrem e com que usuários. \citeonline{frederico2020} observam que, em muitos projetos, o usuário só é consultado depois que o protótipo está pronto, quando alterar o que foi construído é mais caro. Este trabalho parte da perspectiva oposta: avaliar primeiro a usabilidade do aplicativo em uso e, com base nos resultados, orientar as melhorias.

O problema de pesquisa que orienta este trabalho é: \textbf{quais problemas de usabilidade são encontrados no aplicativo de transporte público utilizado em Divinópolis-MG e como esses problemas podem orientar uma proposta de melhoria da experiência do usuário?}

\section{Objetivos do Trabalho}

O objetivo geral deste trabalho é investigar a usabilidade do aplicativo de transporte público utilizado em Divinópolis-MG, identificando dificuldades de interação e propondo melhorias fundamentadas em princípios de usabilidade e experiência do usuário.

Para alcançá-lo, foram definidos os seguintes objetivos específicos:

\begin{alineas}
  \item identificar conceitos, critérios e métodos de avaliação de usabilidade aplicáveis a aplicativos móveis de transporte público;
  \item caracterizar o aplicativo FaleBus, considerando suas funcionalidades e os principais fluxos de interação do usuário;
  \item avaliar a usabilidade do aplicativo com usuários, a partir dos critérios definidos na fundamentação teórica;
  \item identificar e organizar os problemas de usabilidade encontrados, relacionando-os a princípios de usabilidade;
  \item elaborar uma proposta de melhorias de interface e interação fundamentada nos resultados da avaliação.
\end{alineas}

\completar{validar os objetivos específicos com a orientadora, em especial o último, que depende de a proposta incluir ou não um protótipo}

\section{Justificativa}

A avaliação da usabilidade permite identificar problemas a partir da experiência de uso, em vez de presumi-los. \citeonline{suzianti2020} ilustram a relevância desse tipo de avaliação em outro contexto: em um estudo preliminar com 100 usuários do aplicativo da MRT Jakarta, 72\% relataram dificuldades ao usá-lo. Esses números não se aplicam ao FaleBus, mas mostram que problemas de interface podem ser frequentes mesmo em aplicativos oficiais de transporte.

A literatura oferece métricas e métodos consolidados para esse tipo de avaliação \cite{hussain2017}, e há exemplos de estudos que usaram os resultados dos testes para orientar o redesenho da interface \cite{suzianti2020}. Entretanto, os trabalhos analisados no Capítulo \ref{capitulo2} não tratam da usabilidade de um aplicativo de ônibus já em uso em uma cidade de médio porte. Este trabalho contribui ao aplicar esses métodos a um contexto local.

A relevância social do estudo está em seu potencial de contribuir para que o acesso às informações e aos serviços do transporte público seja mais claro para os passageiros. Essa contribuição depende dos resultados da avaliação e não pode ser afirmada de antemão. \completar{incluir dados verificáveis sobre o transporte coletivo de Divinópolis e a base de usuários do aplicativo}

\section{Estrutura da Monografia}

A presente monografia está dividida em sete capítulos. Este primeiro capítulo apresenta o contexto, o problema de pesquisa, os objetivos e a justificativa do trabalho. No Capítulo \ref{capitulo2} são apresentados os trabalhos relacionados e as lacunas que motivam esta pesquisa. No Capítulo \ref{capitulo3} são apresentados os conceitos de usabilidade e experiência do usuário, os métodos de avaliação e os instrumentos que fundamentam o trabalho. No Capítulo \ref{capitulo4} são apresentados os procedimentos metodológicos. No Capítulo \ref{capitulo5} é apresentado o desenvolvimento do trabalho. No Capítulo \ref{capitulo6} são apresentados os resultados e as análises dos mesmos. Por fim, no Capítulo \ref{capitulo7} são apresentadas as conclusões e os possíveis trabalhos futuros.
